\documentclass[12pt]{article}
\usepackage{amsmath}
\usepackage{amsfonts}
\usepackage{float}
\usepackage{fancyhdr}
\usepackage{graphicx}
\usepackage{setspace}
\usepackage[colorlinks=true,linkcolor=blue, citecolor=red]{hyperref}
\usepackage{url}
\usepackage[top=.75in, left=.75in, right=.75in, bottom=1in]{geometry}

% For algorithm
\usepackage{algorithm}
\usepackage{algpseudocode}

% ============ CODE ============
\usepackage{listings}
\usepackage{xcolor}
\definecolor{codegreen}{rgb}{0,0.6,0}
\definecolor{codegray}{rgb}{0.5,0.5,0.5}
\definecolor{codepurple}{rgb}{0.58,0,0.82}
\definecolor{backcolour}{rgb}{0.95,0.95,0.92}

% Styling for the code.
\lstdefinestyle{mystyle}{
    backgroundcolor=\color{backcolour},   
    commentstyle=\color{codegreen},
    keywordstyle=\color{magenta},
    numberstyle=\tiny\color{codegray},
    stringstyle=\color{codepurple},
    basicstyle=\ttfamily\footnotesize,
    breakatwhitespace=false,         
    breaklines=true,                 
    captionpos=b,                    
    keepspaces=true,                 
    numbers=left,                    
    numbersep=5pt,                  
    showspaces=false,                
    showstringspaces=false,
    showtabs=false,                  
    tabsize=2
}
\lstset{style=mystyle}

% Table packages and Unicode / Font support
\usepackage{array}
\usepackage{longtable}
\usepackage{enumitem}
\usepackage{booktabs}
\usepackage[utf8]{inputenc}
\DeclareUnicodeCharacter{2264}{\ensuremath{\le}}
\DeclareUnicodeCharacter{2194}{\ensuremath{\leftrightarrow}}
\usepackage[T1]{fontenc}
\usepackage{lmodern}
\graphicspath{{../../../../templates/}{../../prototype/}{figures/}{./}}
\setlength{\emergencystretch}{3em}
\setlist[enumerate]{label=\arabic*.,leftmargin=*}
\setlist[itemize]{leftmargin=*}

% ============ HEADER AND FOOTER ============
\setlength{\headheight}{29.43912pt}

\pagestyle{fancy}
\fancyfoot{}
\fancyfoot[R]{Page \thepage}

\setlength{\parindent}{0pt}
\renewcommand{\baselinestretch}{1.5}

\newcommand{\coursename}{Introduction of Software Engineering}
\newcommand{\reportname}{Studitopia}
\newcommand{\reporttitle}{\href{https://github.com/RoundFloor/Studitopia.git}{Studitopia}}

\newcommand{\studentname}{%
Nguyen Thanh An (24125024)\\
Nguyen Chi Bao (24125025)\\
Tran Nguyen Khanh Binh (24125026)\\
Tran Canh Anh Tuan (24125048)\\
Nguyen Le Quoc Huy (24125100)
}
\newcommand{\teachername}{Ho Tuan Thanh, Nguyen Thi Minh Tuyen}

% Header
\lhead{\reporttitle}
\rhead{
University of Science - VNUHCM\\
\coursename
}

% Footer
\newcommand{\leftfooter}{}
\lfoot{\leftfooter}

% ============ DOCUMENT ============
\begin{document}

\pagenumbering{roman}
\begin{titlepage}
\newcommand{\HRule}{\rule{\linewidth}{0.5mm}}
\centering

\textsc{\LARGE VNU-HCM}\\[0.5cm]
\textsc{\Large University of Science}\\[0.5cm]
\textsc{\large Faculty of Information Technology}\\[0.5cm]
\textsc{Computer Science}\\[0.5cm]

\HRule \\[0.4cm]
{ 
\huge{\bfseries{\reporttitle}}\\[0.5cm]
\large{\bfseries{UC-31 --- Request Learning Help}}
}\\[0.4cm]
\HRule \\[0.5cm]

\textbf{\large Subject: \coursename}\\[0.5cm]

\begin{minipage}[t]{0.55\textwidth}
\begin{flushleft} \large
\emph{Students:}\\
\studentname
\end{flushleft}
\end{minipage}
~
\begin{minipage}[t]{0.4\textwidth}
\begin{flushright} \large
\emph{Supervising Teachers:} \\
\teachername
\end{flushright}
\end{minipage}\\[1cm]

{\large \today}\\[1cm]

\includegraphics[scale=.20]{img/hcmus-logo.png}\\[1cm] 

\vfill
\end{titlepage}	

\tableofcontents
\pagebreak

\pagenumbering{arabic}
\setcounter{page}{1}

\hypertarget{sec:uc-31-request-learning-help}{}
\addcontentsline{toc}{section}{UC-31 --- Request Learning Help}
\section*{UC-31 --- Request Learning Help}
\par\medskip

\textbf{Version:} \texttt{1.0}
\par\medskip

\textbf{Performed by:} [Tran Nguyen Khanh Binh] | \textbf{Edited by:} [Tran Nguyen Khanh Binh]
\par\medskip

\begingroup
\small
\renewcommand{\baselinestretch}{1.0}\selectfont
\renewcommand{\arraystretch}{1.15}
\setlength{\tabcolsep}{6pt}
\begin{longtable}{|>{\raggedright\arraybackslash}p{0.09\textwidth}|>{\raggedright\arraybackslash}p{0.84\textwidth}|}
\hline
\textbf{Field} & \textbf{Value} \\ \hline
\endhead
Use Case ID & UC-31 \\ \hline
Name & Request Learning Help \\ \hline
Primary Actor & University Student \\ \hline
Supporting Actors & Other Learners (friends, room participants, or concept learners); Notification Service \\ \hline
Feature Group & Supporting Meeting Rooms (Secondary) \\ \hline
Priority & Not explicitly stated \\ \hline
Related User Stories & US-37 \\ \hline
Related Requirements & FR-37.1, FR-37.2, FR-37.3, FR-37.4, FR-37.5, FR-37.6; AC-37.1, AC-37.2, AC-37.3, AC-37.4, AC-37.5, AC-37.6; FR-14.1, FR-14.2, FR-14.3, FR-14.4, FR-14.5, FR-14.6; AC-14.1, AC-14.2, AC-14.3, AC-14.4, AC-14.5; NFR-SEC-5, NFR-SEC-9, NFR-PERF-1, NFR-USE-4, NFR-USE-5 \\ \hline
Source Documents & \texttt{proposals/userstory.md} (10.4); \texttt{proposals/functional\_requirement.md} (US-37); \texttt{proposals/study\_room.md}; \texttt{proposals/hall\_room.md}; \texttt{proposals/meeting\_lobby.md}; \texttt{proposals/settings.md}; \texttt{proposals/non\_functional\_requirement.md}; \texttt{proposals/API.md}; \texttt{proposals/response\_code.md}; \texttt{proposals/feature\_flow.md}; \texttt{proposals/screen\_flow.md} \\ \hline
\end{longtable}
\endgroup
\par\medskip

\hypertarget{sec:part-a-use-case-model}{}
\addcontentsline{toc}{subsection}{Part A --- Use Case Model}
\subsection*{Part A --- Use Case Model}
\par\medskip

\hypertarget{sec:a1-focused-use-case-diagram}{}
\addcontentsline{toc}{subsubsection}{A1. Focused Use Case Diagram}
\subsubsection*{A1. Focused Use Case Diagram}
\par\medskip

\begin{center}
\includegraphics[width=0.85\textwidth,height=0.75\textheight,keepaspectratio]{"figures/UC-31-request-learning-help_1.png"}
\end{center}
\par\medskip

\textbf{Diagram notation.} Person shape = human actor; cloud = supporting system; blue = selected use case; solid line = association; labelled solid arrow = generalization; dashed arrow = include or extend.
\par\medskip

\hypertarget{sec:a2-actor-register}{}
\addcontentsline{toc}{subsubsection}{A2. Actor Register}
\subsubsection*{A2. Actor Register}
\par\medskip

\begingroup
\small
\renewcommand{\baselinestretch}{1.0}\selectfont
\renewcommand{\arraystretch}{1.15}
\setlength{\tabcolsep}{6pt}
\begin{longtable}{|>{\raggedright\arraybackslash}p{0.12\textwidth}|>{\raggedright\arraybackslash}p{0.16\textwidth}|>{\raggedright\arraybackslash}p{0.64\textwidth}|}
\hline
\textbf{Actor} & \textbf{Type} & \textbf{Description} \\ \hline
\endhead
University Student & Primary human actor & The requester who creates, sends, tracks, cancels, or converts a help request. \\ \hline
Other Learners & Secondary human actors & Friends, room participants, or concept learners who receive a help request and can accept or decline it. \\ \hline
Notification Service & Supporting external system & Routes help-request activity to the notification center under \texttt{FR-38.2} / \texttt{AC-38.3}. \\ \hline
\end{longtable}
\endgroup
\par\medskip

\hypertarget{sec:a3-relationships}{}
\addcontentsline{toc}{subsubsection}{A3. Relationships}
\subsubsection*{A3. Relationships}
\par\medskip

\hypertarget{sec:a3-1-associations}{}
\paragraph*{A3.1 Associations}
\par\medskip

\begingroup
\small
\renewcommand{\baselinestretch}{1.0}\selectfont
\renewcommand{\arraystretch}{1.15}
\setlength{\tabcolsep}{6pt}
\begin{longtable}{|>{\raggedright\arraybackslash}p{0.15\textwidth}|>{\raggedright\arraybackslash}p{0.20\textwidth}|>{\raggedright\arraybackslash}p{0.10\textwidth}|>{\raggedright\arraybackslash}p{0.45\textwidth}|}
\hline
\textbf{Actor} & \textbf{Use Case} & \textbf{Type} & \textbf{Description} \\ \hline
\endhead
University Student & UC-31 Request Learning Help & Primary & Creates and manages an outgoing help request. \\ \hline
Other Learners & UC-31 Request Learning Help & Secondary & Receive requests and respond with \textbf{Accept} or \textbf{Decline}. \\ \hline
Notification Service & UC-31 Request Learning Help & Supporting & Delivers help-request notifications to recipients. \\ \hline
\end{longtable}
\endgroup
\par\medskip

\hypertarget{sec:a3-2-generalization-and-specialization}{}
\paragraph*{A3.2 Generalization and Specialization}
\par\medskip

\begingroup
\small
\renewcommand{\baselinestretch}{1.0}\selectfont
\renewcommand{\arraystretch}{1.15}
\setlength{\tabcolsep}{6pt}
\begin{longtable}{|>{\raggedright\arraybackslash}p{0.09\textwidth}|>{\raggedright\arraybackslash}p{0.09\textwidth}|>{\raggedright\arraybackslash}p{0.09\textwidth}|>{\raggedright\arraybackslash}p{0.65\textwidth}|}
\hline
\textbf{Child} & \textbf{Parent} & \textbf{Type} & \textbf{Inherited Behavior or Association} \\ \hline
\endhead
None & None & None & No actor or use-case generalization is named for US-37 in the proposals. \\ \hline
\end{longtable}
\endgroup
\par\medskip

\hypertarget{sec:a3-3-include}{}
\paragraph*{A3.3 Include}
\par\medskip

\begingroup
\small
\renewcommand{\baselinestretch}{1.0}\selectfont
\renewcommand{\arraystretch}{1.15}
\setlength{\tabcolsep}{6pt}
\begin{longtable}{|>{\raggedright\arraybackslash}p{0.14\textwidth}|>{\raggedright\arraybackslash}p{0.18\textwidth}|>{\raggedright\arraybackslash}p{0.60\textwidth}|}
\hline
\textbf{Base Use Case} & \textbf{Included Use Case} & \textbf{Meaning} \\ \hline
\endhead
None & None & No included use case is named for US-37 in the proposals. \\ \hline
\end{longtable}
\endgroup
\par\medskip

\hypertarget{sec:a3-4-extend}{}
\paragraph*{A3.4 Extend}
\par\medskip

\begingroup
\small
\renewcommand{\baselinestretch}{1.0}\selectfont
\renewcommand{\arraystretch}{1.15}
\setlength{\tabcolsep}{6pt}
\begin{longtable}{|>{\raggedright\arraybackslash}p{0.17\textwidth}|>{\raggedright\arraybackslash}p{0.12\textwidth}|>{\raggedright\arraybackslash}p{0.10\textwidth}|>{\raggedright\arraybackslash}p{0.53\textwidth}|}
\hline
\textbf{Extending Use Case} & \textbf{Base Use Case} & \textbf{Condition} & \textbf{Added Behavior} \\ \hline
\endhead
None & None & None & No extending use case is named for US-37 in the proposals. \\ \hline
\end{longtable}
\endgroup
\par\medskip

\hypertarget{sec:a4-model-notes}{}
\addcontentsline{toc}{subsubsection}{A4. Model Notes}
\subsubsection*{A4. Model Notes}
\par\medskip

\begin{enumerate}
\item Only UC-31 and its source-backed actors are shown; no new I-xx, E-xx, or G-xx identifiers are introduced.
\item Help requests are described as generic room help in \texttt{study\_room.md} and \texttt{hall\_room.md}; Hall Room has no separate FR or AC.
\item Forum conversion is part of UC-31 because \texttt{FR-37.5} / \texttt{AC-37.5} require it; the forum remains the primary asynchronous collaboration space.
\end{enumerate}
\par\medskip

\hypertarget{sec:part-b-use-case-specification}{}
\addcontentsline{toc}{subsection}{Part B --- Use Case Specification}
\subsection*{Part B --- Use Case Specification}
\par\medskip

\hypertarget{sec:b1-brief-description}{}
\addcontentsline{toc}{subsubsection}{B1. Brief Description}
\subsubsection*{B1. Brief Description}
\par\medskip

A university student asks for academic assistance from friends, room participants, or concept learners by completing a help form with a title, subject, description, and attachments. \texttt{AC-37.1} identifies the help content, and \texttt{AC-37.2} / \texttt{FR-37.2} identify the permitted recipient groups. Other learners can accept or decline the request; the requester can cancel it or convert it into a forum post for asynchronous discussion (\texttt{AC-37.3}--\texttt{AC-37.5}). The system rate-limits requests to prevent spam (\texttt{AC-37.6}, \texttt{FR-37.6}).
\par\medskip

\hypertarget{sec:b2-preconditions}{}
\addcontentsline{toc}{subsubsection}{B2. Preconditions}
\subsubsection*{B2. Preconditions}
\par\medskip

\begin{enumerate}
\item The student is authenticated with an active session.
\item The student can reach the help area of a room session through \texttt{HelpRequestPanel} / \texttt{CreateHelpRequestButton}.
\item At least one friend, room participant, or concept learner is available as a recipient.
\item The student has not exceeded the help-request rate limit defined by the backend.
\end{enumerate}
\par\medskip

\hypertarget{sec:b3-postconditions}{}
\addcontentsline{toc}{subsubsection}{B3. Postconditions}
\subsubsection*{B3. Postconditions}
\par\medskip

\begin{enumerate}
\item \textbf{Success:} A help request is stored with \texttt{helpTitle}, \texttt{helpSubject}, \texttt{helpDescription}, \texttt{attachments}, and \texttt{recipients} as \texttt{Pending}; recipients receive it and respond with \texttt{Accepted} or \texttt{Declined}.
\item \textbf{Success:} The requester can cancel the request (\texttt{Cancelled}) or convert it into a forum post (\texttt{Converted}).
\item \textbf{Failure:} No new help request is created, and the rate-limit or validation message is shown to the requester.
\end{enumerate}
\par\medskip

\hypertarget{sec:b4-main-success-scenario-basic-flow}{}
\addcontentsline{toc}{subsubsection}{B4. Main Success Scenario (Basic Flow)}
\subsubsection*{B4. Main Success Scenario (Basic Flow)}
\par\medskip

\begin{enumerate}
\item The student selects \textbf{Request Help} in the room's help area.
\item The system opens \texttt{HelpRequestForm}.
\item The student enters \texttt{helpTitle}, \texttt{helpSubject}, and \texttt{helpDescription}, and adds \texttt{attachments}.
\item The student selects \texttt{recipients} from friends, room participants, or concept learners.
\item The student selects \textbf{Send request}.
\item The frontend validates the form fields and attachment file type/size, and the backend checks authentication, recipient eligibility, and the rate limit.
\item The system creates the help request with status \texttt{Pending} and adds it to outgoing help requests.
\item The Notification Service routes the help request to the selected recipients.
\item A recipient reviews the request in \texttt{IncomingHelpRequestList} and selects \textbf{Accept} or \textbf{Decline}.
\item The requester sees the updated status in the outgoing help request list.
\end{enumerate}
\par\medskip

\hypertarget{sec:b5-alternative-flows}{}
\addcontentsline{toc}{subsubsection}{B5. Alternative Flows}
\subsubsection*{B5. Alternative Flows}
\par\medskip

Unless otherwise stated, an alternative flow resumes at the next step of the basic flow.
\par\medskip

\hypertarget{sec:a1-recipient-accepts}{}
\paragraph*{A1. Recipient accepts}
\par\medskip

At basic flow step 10, the recipient selects \textbf{Accept}. The help request status becomes \texttt{Accepted}, and the requester is notified of the response (\texttt{FR-37.3}, \texttt{AC-37.3}).
\par\medskip

\hypertarget{sec:a2-recipient-declines}{}
\paragraph*{A2. Recipient declines}
\par\medskip

At basic flow step 10, the recipient selects \textbf{Decline}. The help request status becomes \texttt{Declined}, and the requester is notified of the response (\texttt{FR-37.3}, \texttt{AC-37.3}).
\par\medskip

\hypertarget{sec:a3-requester-cancels}{}
\paragraph*{A3. Requester cancels}
\par\medskip

The requester selects \textbf{Cancel request} on an outgoing request. The status becomes \texttt{Cancelled}, and no further responses are accepted for that request (\texttt{FR-37.4}, \texttt{AC-37.4}).
\par\medskip

\hypertarget{sec:a4-convert-to-forum-post}{}
\paragraph*{A4. Convert to forum post}
\par\medskip

The requester selects \textbf{Convert to forum post}. The system creates a forum post for asynchronous discussion and records the help request as \texttt{Converted} (\texttt{FR-37.5}, \texttt{AC-37.5}). The forum conversion uses the documented \texttt{POST /forum/posts} contract with permitted post types \texttt{discussion}, \texttt{material}, or \texttt{functional}.
\par\medskip

\hypertarget{sec:b6-exception-flows}{}
\addcontentsline{toc}{subsubsection}{B6. Exception Flows}
\subsubsection*{B6. Exception Flows}
\par\medskip

\hypertarget{sec:e1-rate-limit-reached}{}
\paragraph*{E1. Rate limit reached}
\par\medskip

The backend rejects a new request because the help-request rate limit is reached. The system shows "Too many help requests. Try again later." and does not create the request (\texttt{FR-37.6}, \texttt{AC-37.6}, \texttt{NFR-SEC-9}).
\par\medskip

\hypertarget{sec:e2-required-help-fields-missing-or-attachment-invalid}{}
\paragraph*{E2. Required help fields missing or attachment invalid}
\par\medskip

The system keeps \texttt{HelpRequestForm} open, identifies the invalid title, subject, description, attachment, or recipients, and does not submit the request (\texttt{AC-37.1}, \texttt{NFR-SEC-5}).
\par\medskip

\hypertarget{sec:e3-forum-conversion-fails}{}
\paragraph*{E3. Forum conversion fails}
\par\medskip

If forum post creation fails, the help request is not marked \texttt{Converted}. The documented errors for \texttt{POST /forum/posts} are \texttt{400 INVALID\_POST\_TYPE} (\texttt{48003}), \texttt{400 INVALID\_POST\_CONTENT} (\texttt{48004}), \texttt{401}, \texttt{403 FORBIDDEN} (\texttt{48010}), and \texttt{500 (DATABASE\_ERROR | UNKNOWN\_ERROR)}.
\par\medskip

\hypertarget{sec:b7-business-rules}{}
\addcontentsline{toc}{subsubsection}{B7. Business Rules}
\subsubsection*{B7. Business Rules}
\par\medskip

\begin{enumerate}
\item A help request includes a title, subject, description, and attachments (\texttt{AC-37.1}).
\item Recipients are friends, room participants, or concept learners (\texttt{AC-37.2}, \texttt{FR-37.2}).
\item Other learners can accept or decline a help request (\texttt{AC-37.3}, \texttt{FR-37.3}).
\item The requester can cancel the help request (\texttt{AC-37.4}, \texttt{FR-37.4}).
\item A help request can be converted into a forum post (\texttt{AC-37.5}, \texttt{FR-37.5}).
\item Help requests are rate-limited to prevent spam (\texttt{AC-37.6}, \texttt{FR-37.6}, \texttt{NFR-SEC-9}).
\end{enumerate}
\par\medskip

\hypertarget{sec:b8-special-requirements}{}
\addcontentsline{toc}{subsubsection}{B8. Special Requirements}
\subsubsection*{B8. Special Requirements}
\par\medskip

\begingroup
\small
\renewcommand{\baselinestretch}{1.0}\selectfont
\renewcommand{\arraystretch}{1.15}
\setlength{\tabcolsep}{6pt}
\begin{longtable}{|>{\raggedright\arraybackslash}p{0.10\textwidth}|>{\raggedright\arraybackslash}p{0.83\textwidth}|}
\hline
\textbf{ID} & \textbf{Requirement} \\ \hline
\endhead
NFR-SEC-5 & Uploaded files must be validated for file type and file size before acceptance. \\ \hline
NFR-SEC-9 & The system must rate-limit help requests to reduce spam and abuse. \\ \hline
NFR-PERF-1 & Standard user interactions must respond within 2 seconds. \\ \hline
NFR-USE-4 & Primary interface actions must use clear labels and recognizable controls. \\ \hline
NFR-USE-5 & The user interface must be usable on screen sizes of at least 320 px in width. \\ \hline
\end{longtable}
\endgroup
\par\medskip

\hypertarget{sec:b9-extension-points}{}
\addcontentsline{toc}{subsubsection}{B9. Extension Points}
\subsubsection*{B9. Extension Points}
\par\medskip

\begingroup
\small
\renewcommand{\baselinestretch}{1.0}\selectfont
\renewcommand{\arraystretch}{1.15}
\setlength{\tabcolsep}{6pt}
\begin{longtable}{|>{\raggedright\arraybackslash}p{0.28\textwidth}|>{\raggedright\arraybackslash}p{0.65\textwidth}|}
\hline
\textbf{Extension Point} & \textbf{Location in Flow of Events} \\ \hline
\endhead
Rate-limit check & Basic flow step 6; exception E1 \\ \hline
Accept or decline response & Basic flow step 9; alternatives A1-A2 \\ \hline
Cancellation & After the request is created; alternative A3 \\ \hline
Forum conversion & After the request is created; alternative A4 and exception E3 \\ \hline
\end{longtable}
\endgroup
\par\medskip

\hypertarget{sec:b10-interface-and-data-notes}{}
\addcontentsline{toc}{subsubsection}{B10. Interface and Data Notes}
\subsubsection*{B10. Interface and Data Notes}
\par\medskip

\begin{itemize}
\item \textbf{Screen/flow path:} \texttt{screen\_flow.md} gives \texttt{Home -$>$ Meeting} and \texttt{Meeting -$>$ Hall Room} / \texttt{Meeting -$>$ Study Room}; \texttt{feature\_flow.md} gives the Lobby paths to Hall Room and Study Room. \texttt{study\_room.md} and \texttt{hall\_room.md} place \texttt{HelpRequestPanel} inside the room shell.
\item \textbf{Form fields and actions:} \texttt{HelpRequestForm} collects \texttt{helpTitle} (\textbf{Title}), \texttt{helpSubject} (\textbf{Subject}), \texttt{helpDescription} (\textbf{Description}), and \texttt{attachments} (\textbf{Attachments}). \texttt{HelpRecipientSelector} collects \texttt{recipients} (\textbf{Recipients}) from friends, room participants, or concept learners. Actions are \textbf{Request Help}, \textbf{Send request}, \textbf{Accept}, \textbf{Decline}, \textbf{Cancel request}, and \textbf{Convert to forum post}.
\item \textbf{Statuses:} \texttt{helpStatus} uses \texttt{Pending}, \texttt{Accepted}, \texttt{Declined}, \texttt{Cancelled}, and \texttt{Converted} in outgoing and incoming help lists.
\item \textbf{Exact rate-limit message:} "Too many help requests. Try again later." is the recommended \texttt{HelpRateLimitNotice} copy in \texttt{hall\_room.md}; \texttt{study\_room.md} also shows the shorter "Too many help requests."
\item \textbf{Notification rules:} \texttt{userstory.md} section 11 places Help Request under the Notification Center; \texttt{FR-38.2} / \texttt{AC-38.3} include help requests among social notifications. \texttt{settings.md} exposes \texttt{SocialNotificationToggles}, but no help-specific API is defined.
\item \textbf{Forum conversion API:} \texttt{POST /forum/posts} accepts \texttt{university\_id}, \texttt{major\_id?}, \texttt{course\_id?}, \texttt{post\_type}, \texttt{title}, \texttt{content}, \texttt{tags?}, and \texttt{files?}; success is \texttt{201 \{post\_id, created\_at\}}. Errors and codes are listed in E3 and come from \texttt{API.md} / \texttt{response\_code.md}.
\item \textbf{Help-request endpoints/statuses/error codes:} Not specified in proposals. \texttt{API.md} has no help endpoints, and \texttt{response\_code.md} has no help error codes; only the separate forum-conversion contract is available.
\end{itemize}
\par\medskip

\hypertarget{sec:b11-source-alignment-notes}{}
\addcontentsline{toc}{subsubsection}{B11. Source Alignment Notes}
\subsubsection*{B11. Source Alignment Notes}
\par\medskip

\begin{itemize}
\item UC-07 defines the room session and names the Notification Service for join-request routing; UC-31 uses the same room context for help requests but does not repeat UC-07's join, capacity, or owner-decision rules.
\item UC-05 contributes only the forum-notification pattern after conversion; a converted help request is a new forum post and does not become a comment or answer under UC-05.
\item \texttt{FR-37.1} names title, subject, description, and attachments generally; the exact field names \texttt{helpTitle}, \texttt{helpSubject}, \texttt{helpDescription}, \texttt{attachments}, and \texttt{recipients} come from \texttt{study\_room.md} and \texttt{hall\_room.md}.
\item The Notification Service actor is source-backed by the Notification Center help-request path; \texttt{study\_room.md} and \texttt{hall\_room.md} give slightly different rate-limit copy, and the fuller "Too many help requests. Try again later." is retained.
\end{itemize}
\par\medskip

\par\medskip
\hypertarget{sec:ui-prototype-specimens}{}
\addcontentsline{toc}{subsection}{UI Prototype Specimens}
\subsection*{UI Prototype Specimens}
\par\medskip
\begin{figure}[H]
\centering
\includegraphics[width=0.75\textwidth,height=0.42\textheight,keepaspectratio]{"../../prototype/meetings_e1_request_help.png"}
\par\vspace{2pt}
{\footnotesize \textbf{UI Prototype:} Meetings:E1 Request help}
\end{figure}
\par\smallskip
\begin{figure}[H]
\centering
\includegraphics[width=0.75\textwidth,height=0.42\textheight,keepaspectratio]{"../../prototype/meetings_e2_request_status.png"}
\par\vspace{2pt}
{\footnotesize \textbf{UI Prototype:} Meetings:E2 Request status}
\end{figure}
\par\smallskip
\begin{figure}[H]
\centering
\includegraphics[width=0.75\textwidth,height=0.42\textheight,keepaspectratio]{"../../prototype/meetings_e3_helper_view.png"}
\par\vspace{2pt}
{\footnotesize \textbf{UI Prototype:} Meetings:E3 Helper view}
\end{figure}
\par\smallskip

\end{document}
